\chapter{Common Cold}
\index{Common Cold}

\section*{Definition and Overview}
The common cold, also referred to as acute viral rhinopharyngitis or acute coryza, is a highly contagious, self-limiting viral infection of the upper respiratory tract. It primarily affects the nasal mucosa, sinuses, pharynx, and larynx. Characterised by symptoms such as nasal congestion, rhinorrhoea, sneezing, sore throat, cough, and low-grade fever, it is the most frequent infectious disease in humans. While generally benign, the common cold carries substantial economic impact due to absenteeism from work and school, and it is a major driver of inappropriate antibiotic prescribing.

\section*{Epidemiology}
The common cold occurs globally, with seasonal variations that typically peak during autumn and winter in temperate climates. Adults typically experience two to four episodes per year, whereas young children, particularly those in childcare settings, may experience six to ten episodes annually due to lack of prior immunity and close social contact. Transmission is highly efficient, occurring via direct hand-to-hand contact, self-inoculation of nasal or conjunctival mucosa, and inhalation of airborne respiratory droplets. In many countries, the common cold represents one of the most common reasons for primary care consultations, presenting a significant burden on the healthcare system through diagnostic overlap with more severe respiratory infections.

\section*{Aetiology and Pathophysiology}
A wide variety of respiratory viruses cause the common cold, with host immune responses driving the clinical presentation.
\begin{itemize}
    \item \textbf{Rhinoviruses:} Responsible for 30\% to 80\% of all cases, belonging to the Picornaviridae family, with over 150 known serotypes.
    \item \textbf{Coronaviruses:} The second most common cause, responsible for 10\% to 15\% of cases (excluding SARS-CoV-2).
    \item \textbf{Other viral pathogens:} Influenza viruses, parainfluenza viruses, adenoviruses, respiratory syncytial virus (RSV), enteroviruses, and human metapneumovirus.
    \item \textbf{Viral entry and replication:} The virus binds to cell surface receptors (such as intercellular adhesion molecule 1, or ICAM-1, for most rhinoviruses) on respiratory epithelial cells, replicating locally and initiating an inflammatory response.
    \item \textbf{Inflammatory cascade:} Rather than cellular destruction, the symptoms are primarily caused by the host immune response, including the release of inflammatory mediators (e.g.\ bradykinin, prostaglandins, interleukins, and histamine) that stimulate vascular permeability, mucous secretion, and sensory nerve reflexes.
    \item \textbf{Ciliary dysfunction:} In some infections (like adenovirus), viral replication causes transient damage to the ciliated epithelial cells, impairing mucociliary clearance and predisposing the host to secondary bacterial infections.
\end{itemize}

\section*{Clinical Features}
The clinical features of the common cold typically emerge 1 to 3 days after viral exposure and peak in severity by day 3 or 4.
\begin{itemize}
    \item \textbf{Nasal symptoms:} Rhinorrhoea (initially clear, which may become thick and discoloured without indicating bacterial infection), nasal congestion, and sneezing.
    \item \textbf{Pharyngeal symptoms:} Sore throat, dry throat, or scratchy sensation, which is often the first symptom to appear.
    \item \textbf{Systemic symptoms:} Low-grade fever (more common in infants and young children, rare in adults), mild headache, muscle aches, malaise, and chills.
    \item \textbf{Lower respiratory symptoms:} Mild cough, hoarseness, and post-nasal drip.
    \item \textbf{Signs on examination:} Erythematous nasal mucosa, turbinate swelling, mild pharyngeal erythema, and cervical lymphadenopathy.
\end{itemize}

\section*{Differential Diagnosis}
It is essential to differentiate the common cold from other upper respiratory tract conditions to prevent unnecessary treatment:
\begin{itemize}
    \item \textbf{Allergic rhinitis:} Characterised by pruritus of the nose and eyes, pale/boggy nasal mucosa, sneezing, and absence of sore throat or systemic symptoms.
    \item \textbf{Influenza (flu):} Presents with a sudden onset of high fever, severe myalgia, arthralgia, extreme fatigue, dry cough, and significant headache.
    \item \textbf{Acute bacterial sinusitis:} Suspected when nasal congestion and purulent discharge persist for more than 10 days without improvement, or worsen after initial improvement ("double sickening"), accompanied by facial pain and maxillary tenderness.
    \item \textbf{Streptococcal pharyngitis (strep throat):} Characterised by sudden-onset severe sore throat, high fever, tonsillar exudates, and anterior cervical lymphadenitis, without coryzal symptoms (cough, rhinorrhoea).
    \item \textbf{Pertussis (whooping cough):} Characterised by a prolonged paroxysmal cough followed by an inspiratory "whoop" or post-tussive emesis.
\end{itemize}

\section*{When to Seek Medical Care}
Most cases of the common cold resolve within 7 to 10 days with self-care. Patients should see a general practitioner if symptoms persist past 10 days, worsen suddenly, or if there is persistent high fever.

\textbf{Contact your local emergency services for:}
\begin{itemize}
    \item Difficulty breathing, stridor, or severe shortness of breath.
    \item Cyanosis (blue-coloured lips, tongue, or face).
    \item Chest pain or pressure, particularly when breathing in.
    \item Altered mental state, confusion, or extreme lethargy in infants.
    \item Inability to swallow fluids, drooling, or signs of severe dehydration.
\end{itemize}

\section*{Diagnosis and Investigations}
The diagnosis of the common cold is primarily clinical, based on patient history and physical findings. Investigations are rarely required unless complications are suspected.
\begin{itemize}
    \item \textbf{Clinical evaluation:} History of typical symptom progression and physical examination findings of upper respiratory tract inflammation.
    \item \textbf{Respiratory PCR panel:} Multiplex polymerase chain reaction testing on a nasopharyngeal swab may be performed in hospitalised, immunocompromised, or elderly patients to identify specific viral pathogens (such as influenza, RSV, or SARS-CoV-2).
    \item \textbf{Full blood count:} Not routinely indicated; may show normal or slightly elevated white cell count with a lymphocytic predominance.
    \item \textbf{Sinus imaging:} Plain film X-ray or CT scans are not recommended unless recurrent or chronic bacterial sinusitis is suspected.
\end{itemize}

\section*{Management}
Management focuses on symptomatic relief, patient education, and preventing the misuse of antibiotics.
\begin{itemize}
    \item \textbf{Symptomatic pharmacotherapy:} Oral analgesics (such as paracetamol or ibuprofen) to relieve fever, headache, and myalgias. Avoid aspirin in children due to the risk of Reye's syndrome.
    \item \textbf{Nasal decongestants:} Short-term use (maximum 3 to 5 days) of topical intranasal decongestants (e.g.\ oxymetazoline) or oral decongestants (e.g.\ pseudoephedrine) to relieve nasal obstruction in adults.
    \item \textbf{Saline nasal sprays:} Hypertonic or isotonic saline rinses to moisten nasal passages and thin mucous secretions.
    \item \textbf{Antihistamines:} First-generation antihistamines (e.g.\ chlorpheniramine) combined with decongestants may provide modest reduction in sneezing and rhinorrhoea, though they cause drowsiness.
    \item \textbf{Cough suppressants:} Demulcents (e.g.\ honey for children over 12 months) or cough suppressants (e.g.\ dextromethorphan) for a disruptive dry cough.
    \item \textbf{Hydration and rest:} Adequate fluid intake to maintain hydration and sufficient rest to support immune function.
    \item \textbf{Antibiotic stewardship:} Educating patients that antibiotics are ineffective against viral colds and can cause adverse effects like diarrhoea, rashes, and resistance.
\end{itemize}

\section*{Complications}
\begin{itemize}
    \item \textbf{Acute otitis media:} A common complication in children, occurring when the Eustachian tube is blocked, allowing viral or secondary bacterial infection in the middle ear.
    \item \textbf{Acute rhinosinusitis:} Secondary bacterial infection of the paranasal sinuses, characterised by prolonged purulent nasal discharge and facial pain.
    \item \textbf{Asthma or COPD exacerbation:} Colds are a major trigger for bronchospasm and acute exacerbations of chronic airway diseases.
    \item \textbf{Lower respiratory tract infections:} Progression to viral or secondary bacterial bronchitis, bronchiolitis, or pneumonia, particularly in young children, elderly, or immunocompromised individuals.
    \item \textbf{Rhinitis medicamentosa:} Rebound nasal congestion caused by overuse of topical nasal decongestants beyond 5 consecutive days.
\end{itemize}

\section*{Prognosis and Outcomes}
The prognosis for the common cold is excellent. Most individuals make a full recovery within 7 to 10 days. The cough may persist for up to 3 weeks due to transient airway hyperresponsiveness. Complications are generally mild and respond well to targeted medical interventions. Chronic sequelae are virtually non-existent in healthy hosts, though respiratory symptoms can be prolonged or severe in patients with pre-existing lung conditions.

\section*{Prevention and Self-Care}
\begin{itemize}
    \item \textbf{Frequent hand hygiene:} Washing hands regularly with soap and water for at least 20 seconds, or using alcohol-based hand sanitiser to reduce viral transmission.
    \item \textbf{Respiratory hygiene:} Covering the mouth and nose with a tissue or the inner elbow when coughing or sneezing, and disposing of tissues immediately.
    \item \textbf{Avoid face touching:} Minimising contact with the eyes, nose, and mouth to prevent self-inoculation of viral particles from contaminated surfaces.
    \item \textbf{Surface disinfection:} Cleaning frequently touched surfaces (e.g.\ doorknobs, phones, countertops) regularly, especially when a household member is ill.
    \item \textbf{Social distancing:} Avoiding close contact (such as shaking hands, sharing utensils, or hugging) with individuals showing signs of a cold, and staying home when sick.
\end{itemize}

\section*{Sources and Further Reading}
The following reputable organisations provide further information on the common cold:
\begin{itemize}
    \item Healthdirect Australia. \textit{Understanding the common cold: symptoms, treatment, and prevention.} https://www.healthdirect.gov.au/
    \item Better Health Channel. \textit{Colds (common cold) factsheet.} https://www.betterhealth.vic.gov.au/
    \item Mayo Clinic. \textit{Common cold symptoms, causes, and self-care.} https://www.mayoclinic.org/
    \item National Institute for Health and Care Excellence (NICE). \textit{Prescribing guidelines for respiratory tract infections.} https://www.nice.org.uk/
    \item World Health Organization. \textit{Information on acute respiratory infections.} https://www.who.int/
\end{itemize}
